\scanpage{11}
% archive sheet number: 10
\mspage{- IV-\handwritten{964} -}

\uline{Pour que} $A$ \uline{soit unibranche} (resp. \uline{géométriquement unibranche) , il faut et il suffit que} $\operatorname{Spec}(A)$ \uline{soit irréductible} .

\uline{Proposition} (18.8.1\deleted{3}\handwritten{\unsure{6|5}}) .- \uline{Soit} $A$ \uline{un anneau local noethérien . Si} $A$ \uline{est universellement caténaire , il en est de même de} $ {}^{\mathrm{hs}}A$ .

La démonstration est la même que celle de (18.7.5) .

\uline{18.9. Fibres formelles des anneaux noethériens henséliens} .

\uline{Théorème} (18.9.1) .- \uline{Soit $A$ un anneau local noethérien hensélien , dont les fibres formelles sont géométriquement normales . Alors les fibres formelles de} $A$ \uline{sont géométriquement intègres} (donc \uline{géométriquement connexes}) ,

Les deux assertions de l'énoncé sont en fait équivalentes,un préschéma \inserted{localement noethérien} connexe et normal étant intègre . Comme toute $A$-algèbre intègre finie est un anneau local hensélien (18.5.9 et 18.5.10) , il résulte de (7.3.16.2) que l'on est ramené à prouver que si l'on suppose de plus $A$ \uline{intègre} , alors la fibre du morphisme $\operatorname{Spec}(\hat A)\handwritten{\to}\operatorname{Spec}(A)$ au point générique de $\operatorname{Spec}(A)$ est \uline{intègre} , ce qui résulte du

\uline{Corollaire} (18.9.2) .- \uline{Soit} $A$ \uline{un anneau local noethérien hensélien intègre, dont les fibres formelles sont géométriquement normales . Alors} $\hat A$ \uline{est intè}\deleted{\unsure{géométr.|intègres}}\uline{gre} .

En effet , il résulte de (18.8.12) que $A$ est unibranche , \handwritten{d}onc \deleted{\illegible{2}} le théorème résulte de l'hypothèse faite sur les fibres formelles de $A$ et de (7.6.3) .

\uline{Corollaire} (18.\deleted{\unsure{8}}\handwritten{9}.3) .- \uline{Sous les hypothèses de} (18.9.2) , \uline{le corps des fractions} $L$ \uline{de} $\hat A$ \uline{est une extension séparable du corps des fractions} $K$ \uline{de} $A$ , \uline{e\handwritten{t}} $K$ \uline{est algébriquement fermé dans} $L$ .

\deleted{\illegible{2}} \deleted{La première assertion} \inserted{\handwritten{Cela}} résulte de ce que la fibre du \unsure{morphisme|mbrphisme} $\operatorname{Spec}(\hat A)\handwritten{\to}\handwritten{\to}\operatorname{Spec}(A)$ au point générique de $\operatorname{Spec}(A)$ est géométriquement \deleted{\unsure{intègre}} intègre et de (\handwritten{4.3.2} ) \handwritten{et (4.3.5).}
